% 教材：Shankar, Principles of Quantum Mechanics, Second Edition
% 已覆盖：第 5 章第 5.1--5.5 节，印刷页 pp. 151--176，PDF pp. 163--188
% 人工核验：已完成讲解、教材练习 5.1.1--5.4.3 及第 5.5 节自编应用题

\clearpage
\section{第 5.1 节：自由粒子}
\label{section:QM-C05-S01}

\studymeta
  {QM-C05-S01}
  {2026-09-12}
  {Shankar, \textit{Principles of Quantum Mechanics}, Second Edition}
  {第 5 章 \S 5.1；印刷页 pp.~151--157；PDF pp.~163--169}
  {讲解、教材练习 5.1.1--5.1.4 与订正已完成}

\textbf{本节主线。}

自由粒子没有势能，Hamiltonian 只有动能项 $H=P^2/(2m)$。这一看似最简单的系统第一次把连续谱、简并、广义本征态、传播核、波包扩散和无界算符定义域放在同一套动力学中。核心不是记住若干公式，而是看清三层结构：动量基使演化对角化；位置基中的传播核把同一演化写成积分变换；真正可归一化的粒子态必须由平面波连续叠加成波包。

\textbf{来源边界。}自由粒子谱、传播核、Gaussian 波包和习题 5.1.1--5.1.4 来自教材第 5.1 节。能量归一本征态、Gaussian 宽度记号换算，以及用谱定理区分 $H$ 与 $U(t)$ 的定义域，是对话中为补足严格性而加入的说明。

\notetopic{QM-C05-S01-T01}{自由粒子的能谱、简并与广义本征态}

\keyidea{在动量表象中，自由 Hamiltonian 是乘法算符。能量只依赖 $p^2$，因此每个 $E>0$ 对应相反的两个动量方向；平面波是广义本征态，正规物理态则是它们的连续叠加。}

一维自由粒子的 Hamiltonian 为
\begin{equation}
  H=\frac{P^2}{2m}.
  \label{eq:QM-C05-S01-free-hamiltonian}
\end{equation}
由于 $H$ 是 $P$ 的函数，动量本征态同时是能量本征态：
\begin{equation}
  P\ket{p}=p\ket{p}
  \quad\Longrightarrow\quad
  H\ket{p}=\frac{p^2}{2m}\ket{p}.
  \label{eq:QM-C05-S01-momentum-energy}
\end{equation}
故自由粒子的能谱满足
\begin{equation}
  E=\frac{p^2}{2m}\geq 0.
  \label{eq:QM-C05-S01-energy-spectrum}
\end{equation}
这里的非负性也可直接由
\begin{equation}
  \bra{\psi}H\ket{\psi}
  =\frac{1}{2m}\norm{P\ket{\psi}}^2\geq 0
  \label{eq:QM-C05-S01-positive-hamiltonian}
\end{equation}
看出。

对每个 $E>0$，方程 $p^2/(2m)=E$ 有两支解
\begin{equation}
  p_\alpha(E)=\alpha\sqrt{2mE},
  \qquad \alpha=\pm.
  \label{eq:QM-C05-S01-energy-branches}
\end{equation}
所以 $E>0$ 的能量本征空间为二重简并，可用向右和向左传播的两支标记：
\begin{equation}
  \ket{E,+}\equiv\ket{+\sqrt{2mE}},
  \qquad
  \ket{E,-}\equiv\ket{-\sqrt{2mE}}.
  \label{eq:QM-C05-S01-energy-branch-kets}
\end{equation}
$E=0$ 是谱的端点，两支在 $p=0$ 合并，不能机械地把二重简并结论延伸到该点。

在位置表象中，采用
\begin{equation}
  \braket{x}{p}=\frac{1}{\sqrt{2\pi\hbar}}
  \exp\left(\frac{\ii px}{\hbar}\right),
  \label{eq:QM-C05-S01-plane-wave}
\end{equation}
则固定 $E>0$ 的一般能量本征函数为
\begin{equation}
  \psi_E(x)=A\exp\left(\frac{\ii px}{\hbar}\right)
  +B\exp\left(-\frac{\ii px}{\hbar}\right),
  \qquad p=\sqrt{2mE}.
  \label{eq:QM-C05-S01-energy-general-solution}
\end{equation}
两个系数分别是能量相同、传播方向相反的分量振幅。单个平面波在整条实轴上不平方可积，而满足
\begin{equation}
  \braket{p}{p'}=\delta(p-p'),
  \qquad
  \int_{-\infty}^{\infty}\ket{p}\bra{p}\dd p=I,
  \label{eq:QM-C05-S01-momentum-delta-normalization}
\end{equation}
所以它是连续谱的广义本征态，不是可直接制备的正规化波包。

自由演化在动量基中完全对角：
\begin{equation}
  U(t)=\exp\left(-\frac{\ii P^2t}{2m\hbar}\right)
  =\int_{-\infty}^{\infty}
  \ket{p}\bra{p}\exp\left(-\frac{\ii p^2t}{2m\hbar}\right)\dd p.
  \label{eq:QM-C05-S01-momentum-propagator}
\end{equation}
每个动量分量只获得模长为一的相位，故动量概率分布保持不变。把积分变量由 $p$ 换为 $E$ 时，必须保留两支和 Jacobian，详见
\relatedexercise{QM-C05-S01-E01}{教材练习 5.1.1}。

\dialoguesupplement{若希望把能量本征态本身归一化为 $\delta(E-E')$，应重新定义
\[
  \ket{E,\alpha}_{N}
  =\left(\frac{m}{\sqrt{2mE}}\right)^{1/2}
  \ket{p_\alpha(E)}.
\]
于是 ${}_{N}\!\braket{E,\alpha}{E',\alpha'}_{N}
=\delta_{\alpha\alpha'}\delta(E-E')$。教材习题中的 $\ket{E,\alpha}$ 只是对动量本征态重新标记，因而积分测度中仍显式保留 Jacobian。}

\notetopic{QM-C05-S01-T02}{位置表象中的自由传播核}

\keyidea{传播核 $K(x,t;x',0)=\bra{x}U(t)\ket{x'}$ 是演化算符在位置基中的矩阵元；它把初始位置波函数变换为任意时刻的波函数。}

在位置基中插入动量完备关系，得到
\begin{align}
  K(x,t;x',0)
  &=\bra{x}U(t)\ket{x'}
  \notag\\
  &=\int_{-\infty}^{\infty}
  \braket{x}{p}
  \exp\left(-\frac{\ii p^2t}{2m\hbar}\right)
  \braket{p}{x'}\dd p
  \notag\\
  &=\frac{1}{2\pi\hbar}
  \int_{-\infty}^{\infty}
  \exp\left[
    \frac{\ii p(x-x')}{\hbar}
    -\frac{\ii p^2t}{2m\hbar}
  \right]\dd p.
  \label{eq:QM-C05-S01-kernel-momentum-integral}
\end{align}
令 $d=x-x'$，配方为
\begin{equation}
  pd-\frac{p^2t}{2m}
  =-\frac{t}{2m}\left(p-\frac{md}{t}\right)^2
  +\frac{md^2}{2t}.
  \label{eq:QM-C05-S01-kernel-complete-square}
\end{equation}
利用带有适当收敛规定的 Fresnel Gaussian 积分，得到
\begin{equation}
  \boxed{
  K(x,t;x',0)
  =\sqrt{\frac{m}{2\pi\ii\hbar t}}
  \exp\left[
    \frac{\ii m(x-x')^2}{2\hbar t}
  \right]}
  \qquad(t>0).
  \label{eq:QM-C05-S01-free-kernel}
\end{equation}
指数中的
\begin{equation}
  S_{\mathrm{cl}}(x,t;x',0)
  =\frac{m(x-x')^2}{2t}
  \label{eq:QM-C05-S01-free-classical-action}
\end{equation}
正是连接两个端点的经典自由路径的作用量，因此核的相位可写为 $\exp(\ii S_{\mathrm{cl}}/\hbar)$。

对任意初态，演化为
\begin{equation}
  \boxed{
  \psi(x,t)=\int_{-\infty}^{\infty}
  K(x,t;x',0)\psi(x',0)\dd x'}.
  \label{eq:QM-C05-S01-kernel-evolution}
\end{equation}
这不是先求“粒子从某一点走哪条轨迹”，而是把每个初始位置分量到终点 $x$ 的概率振幅相干相加。由于 $K$ 的量纲是 $L^{-1}$，与 $dd x'$ 相乘后仍给出量纲为 $L^{-1/2}$ 的波函数。

核还必须满足三个一致性条件：
\begin{align}
  \lim_{t\to0^+}K(x,t;x',0)&=\delta(x-x'),
  \label{eq:QM-C05-S01-kernel-delta-limit}\\
  K(x,t_2;x',t_0)
  &=\int K(x,t_2;y,t_1)K(y,t_1;x',t_0)\dd y,
  \label{eq:QM-C05-S01-kernel-composition}\\
  \int K^*(x,t;y,0)K(x,t;y',0)\dd x
  &=\delta(y-y').
  \label{eq:QM-C05-S01-kernel-unitarity}
\end{align}
它们分别对应 $U(0)=I$、演化算符的复合律和酉性。式~\eqref{eq:QM-C05-S01-kernel-delta-limit} 要按分布理解：核本身不逐点趋于普通函数，而是在积分中恢复初始波函数。

\notetopic{QM-C05-S01-T03}{Gaussian 波包的平移与扩散}

\keyidea{自由粒子的平均位置按经典速度平移，动量分布保持不变；但非零动量不确定度使不同动量分量以不同速度传播，从而令位置波包扩散。}

取教材的初始 Gaussian 波包
\begin{equation}
  \psi(x,0)
  =\frac{1}{(\pi\Delta^2)^{1/4}}
  \exp\left(-\frac{x^2}{2\Delta^2}\right)
  \exp\left(\frac{\ii p_0x}{\hbar}\right).
  \label{eq:QM-C05-S01-initial-gaussian}
\end{equation}
这里 $\Delta$ 是指数中的宽度参数，不是位置标准差。直接计算得
\begin{equation}
  \expect{X}_0=0,
  \qquad
  \Delta X(0)=\frac{\Delta}{\sqrt2},
  \qquad
  \expect{P}_0=p_0,
  \qquad
  \Delta P=\frac{\hbar}{\sqrt2\Delta},
  \label{eq:QM-C05-S01-gaussian-initial-moments}
\end{equation}
故 $\Delta X(0)\Delta P=\hbar/2$。

将式~\eqref{eq:QM-C05-S01-initial-gaussian} 代入传播核积分并完成 Gaussian 积分。定义
\begin{equation}
  \tau=\frac{\hbar t}{m\Delta^2},
  \qquad
  v_0=\frac{p_0}{m},
  \label{eq:QM-C05-S01-gaussian-tau}
\end{equation}
可得
\begin{align}
  \psi(x,t)
  ={}&\frac{1}{(\pi\Delta^2)^{1/4}}
  \frac{1}{\sqrt{1+\ii\tau}}
  \exp\left[-\frac{(x-v_0t)^2}{2\Delta^2(1+\ii\tau)}\right]
  \notag\\
  &\times
  \exp\left[\frac{\ii}{\hbar}
  \left(p_0x-\frac{p_0^2t}{2m}\right)\right].
  \label{eq:QM-C05-S01-gaussian-evolution}
\end{align}
取模平方后，相位消失：
\begin{equation}
  \boxed{
  \abs{\psi(x,t)}^2
  =\frac{1}{\sqrt\pi\Delta\sqrt{1+\tau^2}}
  \exp\left[-\frac{(x-v_0t)^2}
  {\Delta^2(1+\tau^2)}\right]}.
  \label{eq:QM-C05-S01-gaussian-density}
\end{equation}
因此
\begin{align}
  \expect{X}_t&=\frac{p_0}{m}t,
  &
  \expect{P}_t&=p_0,
  \label{eq:QM-C05-S01-gaussian-centers}\\
  \Delta X(t)
  &=\frac{\Delta}{\sqrt2}
  \sqrt{1+\left(\frac{\hbar t}{m\Delta^2}\right)^2},
  &
  \Delta P(t)&=\frac{\hbar}{\sqrt2\Delta}.
  \label{eq:QM-C05-S01-gaussian-widths}
\end{align}
波包中心遵循经典自由轨迹，但宽度增长。其特征扩散时间为
\begin{equation}
  t_s=\frac{m\Delta^2}{\hbar}.
  \label{eq:QM-C05-S01-spreading-time}
\end{equation}
若改用实际初始标准差 $\sigma=\Delta X(0)=\Delta/\sqrt2$，同一结果写成
\begin{equation}
  \Delta X(t)=\sigma
  \sqrt{1+\left(\frac{\hbar t}{2m\sigma^2}\right)^2},
  \qquad
  t_s=\frac{2m\sigma^2}{\hbar}.
  \label{eq:QM-C05-S01-gaussian-sigma-convention}
\end{equation}
两套公式没有物理差异，只是宽度参数的定义不同。

扩散不需要外力。初态含有宽度 $\Delta P$ 的动量分布，各分量速度 $p/m$ 不同，传播后逐渐拉开。质量越大或初始位置波包越宽，速度分布越窄，扩散越慢。算符形式
\begin{equation}
  X(t)=X(0)+\frac{P(0)}{m}t,
  \qquad P(t)=P(0)
  \label{eq:QM-C05-S01-heisenberg-free-motion}
\end{equation}
把这一点表达得最直接。利用算符幂级数核对 Gaussian 闭式的计算见
\relatedexercise{QM-C05-S01-E03}{教材练习 5.1.3：Gaussian 演化的另一种求法}。

\notetopic{QM-C05-S01-T04}{定态方程的局部性质与算符定义域}

\keyidea{定态 Schrödinger 方程是二阶微分方程，但势的奇异程度决定波函数及其导数如何衔接；同时，精确酉演化有定义并不意味着可以把无界 Hamiltonian 的指数逐项作用在任意 $L^2$ 态上。}

一维定态方程为
\begin{equation}
  -\frac{\hbar^2}{2m}\psi''(x)+V(x)\psi(x)=E\psi(x),
  \label{eq:QM-C05-S01-stationary-equation}
\end{equation}
即
\begin{equation}
  \psi''(x)=-\frac{2m}{\hbar^2}[E-V(x)]\psi(x).
  \label{eq:QM-C05-S01-second-derivative-equation}
\end{equation}
在势足够规则的区间内，给定某点的 $\psi(x_0)$ 与 $\psi'(x_0)$ 就确定唯一局部解。不同势的连接条件如下：
\begin{itemize}[leftmargin=*]
  \item 若 $V$ 有限且连续，则通常 $\psi$、$\psi'$ 和 $\psi''$ 都连续；
  \item 若 $V$ 有有限跳跃，则 $\psi$ 与 $\psi'$ 连续，而 $\psi''$ 可随 $V$ 跳跃；
  \item 若 $V(x)=g\delta(x-a)$，则 $\psi$ 连续，但对方程跨越 $a$ 积分给出
  \begin{equation}
    \boxed{
    \psi'(a^+)-\psi'(a^-)
    =\frac{2mg}{\hbar^2}\psi(a)}.
    \label{eq:QM-C05-S01-delta-derivative-jump}
  \end{equation}
\end{itemize}
更强的奇异势不能仅凭上述规则猜测，必须从方程及算符定义域重新分析。

教材练习 5.1.4 还揭示了形式级数的陷阱。设自由 Hamiltonian $H=P^2/(2m)$。它是无界算符，在 $\hilbert=L^2(\real)$ 上只定义于真子集
\begin{equation}
  \mathcal D(H)\subsetneq\hilbert.
  \label{eq:QM-C05-S01-h-domain}
\end{equation}
若 $\psi\notin\mathcal D(H)$，则 $H\psi$ 没有 Hilbert 空间中的意义，更不能直接假定
\begin{equation}
  U(t)\psi
  \stackrel{?}{=}
  \sum_{n=0}^{\infty}
  \frac{1}{n!}\left(-\frac{\ii t}{\hbar}\right)^nH^n\psi
  \label{eq:QM-C05-S01-invalid-power-series}
\end{equation}
逐项成立。

但当 $H$ 自伴时，谱定理用有界函数
\begin{equation}
  f_t(E)=\exp\left(-\frac{\ii Et}{\hbar}\right),
  \qquad \abs{f_t(E)}=1
  \label{eq:QM-C05-S01-bounded-spectral-function}
\end{equation}
定义
\begin{equation}
  U(t)=f_t(H)=\exp\left(-\frac{\ii Ht}{\hbar}\right).
  \label{eq:QM-C05-S01-spectral-unitary}
\end{equation}
$U(t)$ 是有界酉算符，且
\begin{equation}
  \boxed{\mathcal D(U(t))=\hilbert}.
  \label{eq:QM-C05-S01-u-domain}
\end{equation}
所以任何 $L^2$ 初态都能精确演化；只是只有当 $\psi\in\mathcal D(H)$ 时，$t=0$ 处的强导数才满足
\begin{equation}
  \left.\frac{\dd}{\dd t}U(t)\psi\right|_{t=0}
  =-\frac{\ii}{\hbar}H\psi.
  \label{eq:QM-C05-S01-strong-derivative-domain}
\end{equation}
教材反例的完整分析见
\relatedexercise{QM-C05-S01-E04}{教材练习 5.1.4：无界生成元的著名反例}。

\textbf{一分钟回顾。}

自由 Hamiltonian 在动量基中对角化，$E=p^2/(2m)$ 使每个正能量对应 $\pm p$ 两个方向。位置传播核是同一酉算符的矩阵元，其积分将初态传播到任意时刻。单个平面波不可正规化，Gaussian 波包则是正规动量叠加：中心按经典轨迹移动，动量分布不变，位置宽度因速度分散而增长。定态方程的连接条件由势的奇异程度决定；在无限维空间中还必须区分无界 $H$ 的定义域与有界酉算符 $U(t)$ 的全空间定义域，不能无条件使用形式幂级数。

\clearpage
\section{第 5.2 节：一维势阱中的粒子}
\label{section:QM-C05-S02}

\studymeta
  {QM-C05-S02}
  {2026-09-13}
  {Shankar, \textit{Principles of Quantum Mechanics}, Second Edition}
  {第 5 章 \S 5.2；印刷页 pp.~157--164；PDF pp.~169--176}
  {讲解、教材练习 5.2.1--5.2.6 与订正已完成}

\textbf{本节主线。}

自由粒子允许在整条实轴上传播，动量连续，因而能量也连续。势阱则把定态 Schr\"odinger 方程与边界条件结合起来：只有少数能量能同时满足方程、连接条件和正规化要求，于是出现离散束缚谱。本节先精确求解无限深方势阱，再用有限深方势阱说明指数尾、超越量子化条件和束缚态阈值，最后讨论基态能量、一般态的时间演化和势阱传播核。

\textbf{来源边界。}无限深方势阱、传播核以及练习 5.2.1--5.2.6 来自教材第 5.2 节。有限阱阈值态的 $L^2$ 严格性、Poincar\'e 不等式给出的精确下界，以及“为何自由粒子常用动量基而势阱常用能量基”的说明，是对话中加入的补充。

\notetopic{QM-C05-S02-T01}{无限深方势阱的边界条件与离散谱}

\keyidea{无限高墙把允许波函数限制在有限区间，并在两端施加齐次 Dirichlet 边界条件。微分方程本身允许连续的 $k$，两端边界条件才把 $k$ 筛成离散序列。}

取长度为 $L$、中心位于原点的无限深方势阱：
\begin{equation}
  V(x)=
  \begin{cases}
    0,&\abs x<L/2,\\
    \infty,&\abs x\geq L/2.
  \end{cases}
  \label{eq:QM-C05-S02-infinite-well-potential}
\end{equation}
势阱外波函数为零；阱内满足
\begin{equation}
  -\frac{\hbar^2}{2m}\psi''(x)=E\psi(x),
  \qquad
  \psi\left(-\frac L2\right)=
  \psi\left(\frac L2\right)=0.
  \label{eq:QM-C05-S02-infinite-well-equation}
\end{equation}
令
\begin{equation}
  k=\frac{\sqrt{2mE}}{\hbar},
  \label{eq:QM-C05-S02-k-definition}
\end{equation}
则阱内一般解为
\begin{equation}
  \psi(x)=A\cos kx+B\sin kx.
  \label{eq:QM-C05-S02-general-inside-solution}
\end{equation}
把两个边界条件写成齐次方程组：
\begin{equation}
  \begin{pmatrix}
    \cos(kL/2)&-\sin(kL/2)\\
    \cos(kL/2)& \sin(kL/2)
  \end{pmatrix}
  \begin{pmatrix}A\\B\end{pmatrix}
  =\begin{pmatrix}0\\0\end{pmatrix}.
  \label{eq:QM-C05-S02-boundary-system}
\end{equation}
要存在非零波函数，系数矩阵必须不可逆，因此
\begin{equation}
  \sin(kL)=0
  \quad\Longrightarrow\quad
  k_n=\frac{n\pi}{L},
  \qquad n=1,2,\ldots.
  \label{eq:QM-C05-S02-quantized-k}
\end{equation}
$n=0$ 只给出零函数，负整数则与正整数给出同一物理本征态，故均不另计。归一化后，能量本征函数可写成
\begin{equation}
  \psi_n(x)=
  \begin{cases}
    \displaystyle\sqrt{\frac2L}\cos\frac{n\pi x}{L},
      &n\text{ 为奇数且 }\abs x<L/2,\\[6pt]
    \displaystyle\sqrt{\frac2L}\sin\frac{n\pi x}{L},
      &n\text{ 为偶数且 }\abs x<L/2,\\[6pt]
    0,&\abs x\geq L/2.
  \end{cases}
  \label{eq:QM-C05-S02-centered-eigenfunctions}
\end{equation}
相应能级为
\begin{equation}
  \boxed{E_n=\frac{\hbar^2\pi^2n^2}{2mL^2}},
  \qquad n=1,2,\ldots.
  \label{eq:QM-C05-S02-infinite-well-energies}
\end{equation}
基态为偶函数，随后奇偶宇称交替。若把势阱平移为 $0<x<L$，本征函数变为 $\sqrt{2/L}\sin(n\pi x/L)$，但能谱不变；完整推导见
\relatedexercise{QM-C05-S02-E05}{教材练习 5.2.5：平移后的无限深势阱}。

\begin{figure}[htbp]
  \centering
  \begin{tikzpicture}[>=Latex,scale=0.88]
    \begin{scope}[xshift=-3.8cm]
      \draw[->] (-2.2,0)--(2.35,0) node[right] {$x$};
      \draw[->] (0,-0.25)--(0,2.35) node[above] {$V$};
      \draw[very thick,noteBlue] (-1.25,2.0)--(-1.25,0)--(1.25,0)--(1.25,2.0);
      \draw[very thick,noteBlue,->] (-1.25,2.0)--(-1.25,2.3);
      \draw[very thick,noteBlue,->] (1.25,2.0)--(1.25,2.3);
      \node[below] at (-1.25,0) {$-L/2$};
      \node[below] at (1.25,0) {$L/2$};
      \node at (0,-0.7) {无限深势阱};
    \end{scope}
    \begin{scope}[xshift=3.8cm]
      \draw[->] (-2.2,0)--(2.35,0) node[right] {$x$};
      \draw[->] (0,-0.25)--(0,2.35) node[above] {$V$};
      \draw[very thick,accentRed] (-2.05,1.55)--(-1.25,1.55)--(-1.25,0)--(1.25,0)--(1.25,1.55)--(2.05,1.55);
      \node[left] at (-0.05,1.55) {$V_0$};
      \node[below] at (-1.25,0) {$-a$};
      \node[below] at (1.25,0) {$a$};
      \node at (0,-0.7) {有限深势阱};
    \end{scope}
  \end{tikzpicture}
  \caption{无限深与有限深方势阱。有限高墙允许束缚态波函数进入经典禁区并形成指数衰减尾。}
  \label{fig:QM-C05-S02-well-potentials}
\end{figure}

\textbf{无限墙处的连接条件。}有限势阶跃处通常要求 $\psi$ 与 $\psi'$ 连续，但不能把这条结论直接套到 $V=\infty$ 的奇异理想化上。无限深势阱只要求波函数在墙上为零；阱内导数一般不必与阱外零函数的导数匹配。

\dialoguesupplement{阱内的正弦或余弦可以形式地拆成 $\ee^{\pm\ii kx}$，但这不表示动量测量只会精确得到 $\pm\hbar k$。截断在有限区间内的波函数，其整条实轴 Fourier 变换通常给出连续动量分布，只在 $\pm\hbar k$ 附近出现主要峰。硬墙还破坏平移对称性，使 $P$ 不再与势阱 Hamiltonian 共同对角化；因此 $\hbar k_n$ 在这里首先是波数对应的动量尺度，而不是该定态的确定动量。}

\notetopic{QM-C05-S02-T02}{有限深势阱、指数尾与束缚态量子化}

\keyidea{有限墙不要求波函数在边界消失，而要求阱内振荡解与阱外衰减解平滑连接。正规化删除无穷远处的增长解，边界匹配再把连续能量参数筛成离散根。}

考虑
\begin{equation}
  V(x)=
  \begin{cases}
    0,&\abs x\leq a,\\
    V_0,&\abs x>a.
  \end{cases}
  \label{eq:QM-C05-S02-finite-well-potential}
\end{equation}
束缚态满足 $0<E<V_0$。定义
\begin{equation}
  k=\frac{\sqrt{2mE}}{\hbar},
  \qquad
  \kappa=\frac{\sqrt{2m(V_0-E)}}{\hbar},
  \label{eq:QM-C05-S02-finite-k-kappa}
\end{equation}
则
\begin{equation}
  k^2+\kappa^2=\frac{2mV_0}{\hbar^2}.
  \label{eq:QM-C05-S02-k-kappa-circle}
\end{equation}
阱内解振荡，阱外解指数变化。势关于原点对称，所以可分别取定宇称。对 $x\geq0$，偶态与奇态分别可写成
\begin{align}
  \psi_{\mathrm e}(x)&=
  \begin{cases}
    A\cos kx,&0\leq x\leq a,\\
    C\ee^{-\kappa x},&x>a,
  \end{cases}
  \label{eq:QM-C05-S02-even-piecewise}\\
  \psi_{\mathrm o}(x)&=
  \begin{cases}
    A\sin kx,&0\leq x\leq a,\\
    C\ee^{-\kappa x},&x>a.
  \end{cases}
  \label{eq:QM-C05-S02-odd-piecewise}
\end{align}
左半轴由相应宇称确定。有限阶跃处 $\psi$ 与 $\psi'$ 连续。偶态在 $x=a$ 给出
\begin{equation}
  A\cos ka=C\ee^{-\kappa a},
  \qquad
  -Ak\sin ka=-\kappa C\ee^{-\kappa a}.
  \label{eq:QM-C05-S02-even-matching}
\end{equation}
两式相除消去振幅，得到
\begin{equation}
  \boxed{k\tan ka=\kappa}.
  \label{eq:QM-C05-S02-even-transcendental}
\end{equation}
奇态同理：
\begin{equation}
  A\sin ka=C\ee^{-\kappa a},
  \qquad
  Ak\cos ka=-\kappa C\ee^{-\kappa a},
  \label{eq:QM-C05-S02-odd-matching}
\end{equation}
故
\begin{equation}
  \boxed{k\cot ka=-\kappa}.
  \label{eq:QM-C05-S02-odd-transcendental}
\end{equation}

令
\begin{equation}
  \alpha=ka,
  \qquad
  \beta=\kappa a,
  \qquad
  z_0^2=\frac{2mV_0a^2}{\hbar^2},
  \label{eq:QM-C05-S02-dimensionless-variables}
\end{equation}
则允许束缚态是第一象限中圆
\begin{equation}
  \alpha^2+\beta^2=z_0^2
  \label{eq:QM-C05-S02-dimensionless-circle}
\end{equation}
与曲线
\begin{equation}
  \beta=\alpha\tan\alpha
  \quad\text{或}\quad
  \beta=-\alpha\cot\alpha
  \label{eq:QM-C05-S02-dimensionless-curves}
\end{equation}
的交点。每个交点确定一个离散的 $\alpha$，从而确定
\begin{equation}
  E=\frac{\hbar^2\alpha^2}{2ma^2}.
  \label{eq:QM-C05-S02-energy-from-alpha}
\end{equation}

第一条偶态分支从 $\alpha=0$ 附近的零值连续上升，因此任意 $z_0>0$ 的圆都至少与它相交一次。第一条奇态分支则始于 $\alpha=\pi/2$；真正出现正规化奇束缚态需要
\begin{equation}
  z_0>\frac\pi2
  \quad\Longleftrightarrow\quad
  \boxed{V_0>\frac{\pi^2\hbar^2}{8ma^2}}.
  \label{eq:QM-C05-S02-first-odd-threshold}
\end{equation}
等号处 $\alpha=\pi/2$、$\beta=0$ 且 $E=V_0$。此时阱外解不再衰减，而是 $C+Dx$；除零函数外不能平方可积，所以它只是连续谱边缘的 threshold solution（阈值解），不是真正束缚态。教材把等号包含在“出现条件”中时，应按阈值意义理解。完整图解与极限分析见
\relatedexercise{QM-C05-S02-E06}{教材练习 5.2.6：有限深方势阱}。

当 $V_0\to\infty$ 时，对固定低能态有 $\beta\to\infty$。偶态根趋于 $(2r+1)\pi/2$，奇态根趋于 $r\pi$，两族合并为
\begin{equation}
  \alpha_n=\frac{n\pi}{2},
  \qquad
  E_n=\frac{\hbar^2\pi^2n^2}{8ma^2},
  \label{eq:QM-C05-S02-infinite-limit}
\end{equation}
正好恢复宽度 $2a$ 的无限深势阱。

\textbf{离散束缚谱与连续散射谱。}束缚态必须在左右无穷远同时衰减，因而要满足两端正规化条件和所有边界连接条件；一般只有离散能量能同时满足这些约束。若 $E>V_0$，阱外也变为振荡解，粒子可传播至无穷远，渐近波数可连续变化，得到连续散射谱。离散性来自全局边界条件，不只是“势阱有墙”。例如谐振子虽然没有有限位置的墙，但 $V(x)\to\infty$ 仍要求两端衰减，因此同样产生离散能级。

\dialoguesupplement{若把式~\eqref{eq:QM-C05-S02-finite-well-potential} 整体减去 $V_0$，就得到阱内为 $-V_0$、无穷远为零的纯吸引势；波函数和束缚态数目不变，能量整体变为 $E-V_0$。练习 5.2.2 用变分法说明：在适当局域性条件下，一维任意非零纯吸引势至少存在一个束缚态。奇异的吸引 delta 势同样只有一个束缚态，但其导数在原点发生有限跳跃，见练习 5.2.3。}

\notetopic{QM-C05-S02-T03}{基态能量、时间演化与表象选择}

\keyidea{严格局域在有限区间内的非零波函数不可能同时满足 $\Delta P=0$；因此势阱基态具有非零能量。能量本征基使时间演化对角化，而位置表象则把同一离散谱写成传播核。}

无限深势阱不存在 $E=0$ 的非零定态。若 $E=0$，阱内方程给出
\begin{equation}
  \psi''(x)=0
  \quad\Longrightarrow\quad
  \psi(x)=A+Bx.
  \label{eq:QM-C05-S02-zero-energy-linear}
\end{equation}
两端条件 $\psi(\pm L/2)=0$ 迫使 $A=B=0$。因此基态能量严格为正：
\begin{equation}
  E_1=\frac{\pi^2\hbar^2}{2mL^2}>0.
  \label{eq:QM-C05-S02-ground-energy}
\end{equation}

教材用不确定性原理作数量级估计。定态可取为实函数，故 $\expect P=0$，并有
\begin{equation}
  \expect H=\frac{(\Delta P)^2}{2m}.
  \label{eq:QM-C05-S02-energy-from-delta-p}
\end{equation}
由支集限制 $\Delta X\leq L/2$ 和
\begin{equation}
  \Delta X\Delta P\geq\frac\hbar2
  \label{eq:QM-C05-S02-uncertainty-bound}
\end{equation}
可得粗略下界
\begin{equation}
  \expect H\geq\frac{\hbar^2}{2mL^2}.
  \label{eq:QM-C05-S02-rough-ground-bound}
\end{equation}
它正确给出 $\hbar^2/(mL^2)$ 的尺度，却少了精确系数 $\pi^2$。原因是 $\Delta X\leq L/2$ 与 Heisenberg 不等式的等号条件不能在硬墙边界下同时达到，不能把两条不等式都当成等式。

更精确地，对满足 $\psi(\pm L/2)=0$ 的归一化波函数，Poincar\'e 不等式给出
\begin{equation}
  \int_{-L/2}^{L/2}\abs{\psi'(x)}^2\dd x
  \geq
  \left(\frac\pi L\right)^2
  \int_{-L/2}^{L/2}\abs{\psi(x)}^2\dd x.
  \label{eq:QM-C05-S02-poincare-inequality}
\end{equation}
因此
\begin{equation}
  \expect H
  =\frac{\hbar^2}{2m}
  \int_{-L/2}^{L/2}\abs{\psi'(x)}^2\dd x
  \geq\frac{\pi^2\hbar^2}{2mL^2},
  \label{eq:QM-C05-S02-exact-variational-bound}
\end{equation}
且等号由基态余弦取得。一般的变分下界及其对一维吸引势的应用见
\relatedexercise{QM-C05-S02-E02}{教材练习 5.2.2：变分下界与一维吸引势}。

任意初态可在箱中能量本征基展开：
\begin{equation}
  \ket{\psi(0)}=\sum_{n=1}^{\infty}c_n\ket n,
  \qquad
  c_n=\braket{n}{\psi(0)}.
  \label{eq:QM-C05-S02-general-initial-expansion}
\end{equation}
时间演化为
\begin{equation}
  \boxed{
  \ket{\psi(t)}=
  \sum_{n=1}^{\infty}
  c_n\ee^{-\ii E_nt/\hbar}\ket n}.
  \label{eq:QM-C05-S02-general-time-evolution}
\end{equation}
能量测量概率 $\abs{c_n}^2$ 不随时间改变，但位置概率密度一般含有
\begin{equation}
  c_m^*c_n\psi_m^*(x)\psi_n(x)
  \ee^{-\ii(E_n-E_m)t/\hbar}
  \label{eq:QM-C05-S02-position-cross-terms}
\end{equation}
这样的交叉项，所以叠加态的 $\abs{\psi(x,t)}^2$ 通常随时间变化。只有单一能量本征态或同一简并能级内的叠加才是 stationary state（定态）。

传播算符和位置传播核分别为
\begin{align}
  U(t)
  &=\sum_{n=1}^{\infty}
  \ket n\bra n\ee^{-\ii E_nt/\hbar},
  \label{eq:QM-C05-S02-box-propagator-operator}\\
  K(x,t;x',0)
  &=\bra xU(t)\ket{x'}
  =\sum_{n=1}^{\infty}
  \psi_n(x)\psi_n^*(x')
  \ee^{-\ii E_nt/\hbar}.
  \label{eq:QM-C05-S02-box-kernel}
\end{align}
与自由粒子的动量积分相比，这里出现离散求和，是因为箱中 Hamiltonian 的束缚谱离散。

\dialoguesupplement{选择动量基还是能量基，不取决于“动量比能量更基本”，而取决于哪组完备基能把 Hamiltonian 对角化。自由粒子有平移对称性，$[H,P]=0$，所以动量本征态同时是能量本征态，常写成动量积分。硬墙破坏平移对称性，箱中定态一般不是 $P$ 的本征态，直接用能量本征态求和更自然。仍可把箱中波函数 Fourier 变换到连续动量表象，但势能会耦合不同动量，Hamiltonian 不再对角。环形周期边界同时保留离散动量和平移对称性，说明真正判据是“是否对角化 $H$”，而不是“离散就用能量、连续就用动量”。}

势阱参数改变时还需区分两种极限。若墙突然移动，瞬间波函数来不及改变，之后要把旧波函数投影到新能量基；见
\relatedexercise{QM-C05-S02-E01}{教材练习 5.2.1：势阱突然扩张}。若墙缓慢移动且满足绝热条件，粒子保持在对应的瞬时本征态，能量变化给出广义力
\begin{equation}
  F=-\pdv{E_n}{L};
  \label{eq:QM-C05-S02-adiabatic-force}
\end{equation}
其与经典碰壁平均力的一致性见
\relatedexercise{QM-C05-S02-E04}{教材练习 5.2.4：墙壁所受平均力}。

\textbf{一分钟回顾。}

无限深势阱用 $\psi(\pm L/2)=0$ 把连续波数筛为 $k_n=n\pi/L$，从而得到离散、正的能谱。有限墙允许指数尾，但正规化与边界匹配仍只容许离散束缚能量；高于阱外势的平台则形成连续散射谱。基态非零能量来自空间局域化，Poincar\'e 不等式给出精确下界。时间演化在能量基中只是逐分量积累相位，能量概率不变而位置干涉项可随时间变化。表象的选择应服从“让 Hamiltonian 尽量对角”的计算原则。

\clearpage
\section{第 5.3 节：概率连续性方程}
\label{section:QM-C05-S03}

\studymeta
  {QM-C05-S03}
  {2026-09-14}
  {Shankar, \textit{Principles of Quantum Mechanics}, Second Edition}
  {第 5 章 \S 5.3；印刷页 pp.~165--167；PDF pp.~177--179}
  {讲解、教材练习 5.3.1--5.3.4 与订正已完成}

\textbf{本节主线。}

酉演化保证全空间总概率不变，但只知道“总量不变”还不能说明概率如何在空间中转移。连续性方程把整体守恒加强为局域守恒：任一区域内概率的变化只能由穿过边界的概率流解释。对标准的实标势 Hamiltonian，这一方程不是额外公设，而是 Born 概率密度与 Schr\"odinger 方程的直接推论。

\textbf{来源边界。}整体与局域守恒的对比、概率流的推导、系综解释以及教材练习 5.3.1--5.3.4 来自教材第 5.3 节。振幅--相位分解、stationary state 不必满足 $\mathbf j=0$、实波函数不等于零动量本征态，以及复势作为开放通道有效模型的解释，是对话中的严格性与物理直觉补充。

\notetopic{QM-C05-S03-T01}{整体概率守恒与局域连续性方程}

\keyidea{整体概率守恒只约束全空间积分；局域守恒进一步要求，固定区域内概率的减少必须等于穿过其边界的净向外概率流。}

定义三维位置概率密度
\begin{equation}
  \rho(\mathbf r,t)=\abs{\psi(\mathbf r,t)}^2.
  \label{eq:QM-C05-S03-density}
\end{equation}
若封闭系统按酉算符 $U(t)$ 演化，则
\begin{align}
  \int_{\real^3}\rho(\mathbf r,t)\dd^3r
  &=\braket{\psi(t)}{\psi(t)} \\
  &=\bra{\psi(0)}U^\dagger(t)U(t)\ket{\psi(0)}
  =\braket{\psi(0)}{\psi(0)}.
  \label{eq:QM-C05-S03-global-conservation}
\end{align}
归一化初态因而始终满足全空间总概率为一。这是 global conservation（整体守恒），却不排除“某处消失、远处同时出现而总量碰巧不变”的非局域重排。局域守恒要求存在 probability current density（概率流密度）$\mathbf j$，使
\begin{equation}
  \boxed{
  \pdv{\rho}{t}+\nabla\cdot\mathbf j=0}.
  \label{eq:QM-C05-S03-continuity}
\end{equation}

对固定区域 $\Omega$ 定义
\begin{equation}
  P_\Omega(t)=\int_\Omega\rho(\mathbf r,t)\dd^3r.
  \label{eq:QM-C05-S03-region-probability}
\end{equation}
将式~\eqref{eq:QM-C05-S03-continuity} 在 $\Omega$ 上积分，并使用散度定理，得到
\begin{equation}
  \boxed{
  \odv{P_\Omega}{t}
  =-\oint_{\partial\Omega}\mathbf j\cdot\dd\mathbf S}.
  \label{eq:QM-C05-S03-integral-continuity}
\end{equation}
$\dd\mathbf S=\mathbf n\dd S$ 按约定指向区域外侧。边界上的净向外通量为正时，区域内概率减少；净向内通量为正时，区域内概率增加。

\begin{figure}[htbp]
  \centering
  \begin{tikzpicture}[>=Latex,scale=0.9]
    \draw[thick,fill=blue!4]
      plot[smooth cycle,tension=0.8]
      coordinates {(-2.4,0.1) (-1.5,1.45) (0.2,1.7) (2.2,0.8)
                   (2.45,-0.5) (0.8,-1.5) (-1.4,-1.25)};
    \node at (0,0) {$\Omega$};
    \node[above] at (0.2,1.7) {$\partial\Omega$};
    \draw[->,thick,blue!70!black] (-1.85,0.95)--(-2.75,1.55);
    \draw[->,thick,blue!70!black] (1.65,1.2)--(2.25,2.0);
    \draw[->,thick,blue!70!black] (2.35,-0.1)--(3.35,-0.15);
    \draw[->,thick,blue!70!black] (0.2,-1.45)--(0.15,-2.35);
    \node[right,blue!70!black] at (3.35,-0.15)
      {$\mathbf j\cdot\mathbf n$};
    \draw[->,thick,red!70!black] (-3.2,-0.55)--(-2.15,-0.45);
    \node[left,red!70!black] at (-3.2,-0.55) {流入};
  \end{tikzpicture}
  \caption{固定区域内概率的变化由穿过边界的净通量决定。蓝色箭头表示沿外法向的流出，红色箭头表示流入。}
  \label{fig:QM-C05-S03-boundary-flux}
\end{figure}

一维区间 $[a,b]$ 上的积分形式尤其直观：
\begin{equation}
  \boxed{
  \odv{}{t}\int_a^b\rho(x,t)\dd x
  =j(a,t)-j(b,t)}.
  \label{eq:QM-C05-S03-one-dimensional-balance}
\end{equation}
这里约定 $j>0$ 指向 $+x$：$j(a)$ 是从左端进入的有向流，$j(b)$ 是从右端离开的有向流；公式本身也自动处理反向流。

将 $\Omega$ 扩展到全空间，若无穷远边界满足
\begin{equation}
  \lim_{R\to\infty}\oint_{S_R}\mathbf j\cdot\dd\mathbf S=0,
  \label{eq:QM-C05-S03-no-flux-at-infinity}
\end{equation}
则局域连续性方程推出整体概率守恒。反过来，只有总积分不变不能推出局域连续性方程；在无穷远无通量的通常情形下，整体守恒只是局域守恒的必要而非充分条件。

\notetopic{QM-C05-S03-T02}{从 Schr\"odinger 方程推导概率流}

\keyidea{对 Schr\"odinger 方程及其复共轭作适当的乘法与相减，势能项因实标势而抵消，动能项则可由乘积法则写成一个散度。}

对实值标势 $V(\mathbf r,t)$，位置表象中的 Schr\"odinger 方程及其复共轭分别为
\begin{align}
  \ii\hbar\pdv{\psi}{t}
  &=-\frac{\hbar^2}{2m}\nabla^2\psi+V\psi,
  \label{eq:QM-C05-S03-schrodinger}\\
  -\ii\hbar\pdv{\psi^*}{t}
  &=-\frac{\hbar^2}{2m}\nabla^2\psi^*+V\psi^*.
  \label{eq:QM-C05-S03-conjugate-schrodinger}
\end{align}
第二式使用了 $V^*=V$；在适当定义域和边界条件下，这对应乘法势算符的 Hermitian 性。用 $\psi^*$ 乘第一式，用 $\psi$ 乘第二式，再将两式相减：
\begin{align}
  \ii\hbar
  \left(
    \psi^*\pdv{\psi}{t}
    +\psi\pdv{\psi^*}{t}
  \right)
  &=-\frac{\hbar^2}{2m}
  \left(
    \psi^*\nabla^2\psi
    -\psi\nabla^2\psi^*
  \right),\\
  \ii\hbar\pdv{\rho}{t}
  &=-\frac{\hbar^2}{2m}
  \left(
    \psi^*\nabla^2\psi
    -\psi\nabla^2\psi^*
  \right).
  \label{eq:QM-C05-S03-density-time-derivative}
\end{align}
势能项确实消失，因为
\begin{equation}
  \psi^*V\psi-\psi V\psi^*=0.
  \label{eq:QM-C05-S03-potential-cancellation}
\end{equation}

乘积法则给出
\begin{align}
  \nabla\cdot(\psi^*\nabla\psi)
  &=\nabla\psi^*\cdot\nabla\psi+\psi^*\nabla^2\psi,\\
  \nabla\cdot(\psi\nabla\psi^*)
  &=\nabla\psi\cdot\nabla\psi^*+\psi\nabla^2\psi^*.
\end{align}
相减后，两个梯度内积相同而抵消，因此
\begin{equation}
  \psi^*\nabla^2\psi-\psi\nabla^2\psi^*
  =\nabla\cdot
  (\psi^*\nabla\psi-\psi\nabla\psi^*).
  \label{eq:QM-C05-S03-divergence-identity}
\end{equation}
代回式~\eqref{eq:QM-C05-S03-density-time-derivative} 并与连续性方程比较，得到
\begin{equation}
  \boxed{
  \mathbf j
  =\frac{\hbar}{2m\ii}
  (\psi^*\nabla\psi-\psi\nabla\psi^*)
  =\frac{\hbar}{m}\operatorname{Im}(\psi^*\nabla\psi)}.
  \label{eq:QM-C05-S03-probability-current}
\end{equation}
一维时只需把 $\nabla$ 换成 $\partial_x$：
\begin{equation}
  j(x,t)=\frac{\hbar}{2m\ii}
  \left(
    \psi^*\pdv{\psi}{x}
    -\psi\pdv{\psi^*}{x}
  \right).
  \label{eq:QM-C05-S03-current-one-dimensional}
\end{equation}
因此 $\rho$ 来自 Born 规则，而 $\mathbf j$ 的具体形式由 Schr\"odinger 动力学和局域守恒共同确定，并不是另加的独立概率公设。

\notetopic{QM-C05-S03-T03}{相位梯度、定态概率流与源汇}

\keyidea{概率密度由波函数振幅决定，概率流还读取空间相位梯度。定态只要求流的散度为零，不要求流本身为零；非 Hermitian 有效势则会在连续性方程中产生源项或汇项。}

把波函数写成振幅--相位形式
\begin{equation}
  \psi(\mathbf r,t)=R(\mathbf r,t)
  \exp\left(\frac{\ii S(\mathbf r,t)}{\hbar}\right),
  \qquad R\geq0,
  \label{eq:QM-C05-S03-polar-wavefunction}
\end{equation}
则
\begin{equation}
  \rho=R^2,
  \qquad
  \boxed{\mathbf j=\frac{\rho}{m}\nabla S}.
  \label{eq:QM-C05-S03-current-phase}
\end{equation}
所以振幅决定局部概率密度，相位梯度决定概率的局部流向。只要 $\rho\neq0$，还可定义
\begin{equation}
  \mathbf v_{\mathrm{flow}}
  =\frac{\mathbf j}{\rho}
  =\frac{\nabla S}{m}.
  \label{eq:QM-C05-S03-flow-velocity}
\end{equation}
它是概率流的局域速度场，不等于一般叠加态的一次动量测量结果；在波函数节点 $\rho=0$ 处也不应使用该比值。

对动量为 $\mathbf p$ 的平面波，
\begin{equation}
  \psi_{\mathbf p}(\mathbf r,t)
  =A\exp\left[
    \frac{\ii}{\hbar}
    (\mathbf p\cdot\mathbf r-Et)
  \right],
  \label{eq:QM-C05-S03-plane-wave}
\end{equation}
有
\begin{equation}
  \rho=\abs A^2,
  \qquad
  \boxed{\mathbf j=\rho\frac{\mathbf p}{m}}.
  \label{eq:QM-C05-S03-plane-wave-current}
\end{equation}
这里 $\rho$ 和 $\mathbf j$ 都是常量，因此 $\partial_t\rho=0$、$\nabla\cdot\mathbf j=0$。平面波虽不可普通归一化，却仍在形式上满足局域连续性方程；完整计算见
\relatedexercise{QM-C05-S03-E03}{教材练习 5.3.3：动量本征函数的概率流}。

若 $\psi(x,t)=\phi(x)\ee^{-\ii Et/\hbar}$ 是 stationary state，则 $\partial_t\rho=0$，一维连续性方程只推出
\begin{equation}
  \pdv{j}{x}=0
  \quad\Longrightarrow\quad
  j=\text{常数},
  \label{eq:QM-C05-S03-stationary-constant-current}
\end{equation}
而不自动推出 $j=0$。自由平面波就是具有非零恒定流的定态。对通常的一维正规化束缚态，无穷远处没有概率流，常数才被迫为零；若空间函数可取为实函数，式~\eqref{eq:QM-C05-S03-current-one-dimensional} 也直接给出 $j=0$。

\dialoguesupplement{实空间波函数只保证净概率流为零，不保证动量测量必得零。例如 $\cos kx=(\ee^{\ii kx}+\ee^{-\ii kx})/2$ 是 $\pm\hbar k$ 两个动量方向的等权叠加，二者的流抵消，但动量涨落并不为零。相应推导见\relatedexercise{QM-C05-S03-E02}{教材练习 5.3.2：实波函数的概率流}。}

更一般地，一维左右行波叠加
\begin{equation}
  \psi(x)=A\ee^{\ii px/\hbar}+B\ee^{-\ii px/\hbar}
  \label{eq:QM-C05-S03-counterpropagating-state}
\end{equation}
具有概率密度
\begin{equation}
  \rho(x)=\abs A^2+\abs B^2
  +2\operatorname{Re}
  \left(AB^*\ee^{2\ii px/\hbar}\right),
  \label{eq:QM-C05-S03-counterpropagating-density}
\end{equation}
但概率流为
\begin{equation}
  \boxed{j=\frac{p}{m}(\abs A^2-\abs B^2)}.
  \label{eq:QM-C05-S03-counterpropagating-current}
\end{equation}
$\rho$ 中的交叉项形成空间干涉条纹；$j$ 的反对称导数组合则使这些交叉项抵消，只留下右行流减左行流。相对相位会平移条纹，却不改变净流。完整推导见
\relatedexercise{QM-C05-S03-E04}{教材练习 5.3.4}。

概率流还有直接的 ensemble interpretation（系综解释）：若 $N$ 个独立、同样制备的粒子都处于态 $\psi$，则短时间 $\dd t$ 内穿过有向面积元 $\dd\mathbf S$ 的期望粒子数为
\begin{equation}
  \dd N=N\mathbf j\cdot\dd\mathbf S\dd t.
  \label{eq:QM-C05-S03-ensemble-flux}
\end{equation}
三维中 $[\rho]=L^{-3}$、$[\mathbf j]=L^{-2}T^{-1}$，与这一解释一致。

最后，若采用复有效势
\begin{equation}
  V=V_R-\ii V_I,
  \qquad V_I>0,
  \label{eq:QM-C05-S03-complex-potential}
\end{equation}
则 Hamiltonian 不再 Hermitian，连续性方程变为
\begin{equation}
  \boxed{
  \pdv{\rho}{t}+\nabla\cdot\mathbf j
  =-\frac{2V_I}{\hbar}\rho}.
  \label{eq:QM-C05-S03-sink-continuity}
\end{equation}
右端是 probability sink（概率汇）。若 $V_I$ 为常数、无穷远没有通量且初始总概率为一，则
\begin{equation}
  P(t)=\exp\left(-\frac{2V_It}{\hbar}\right).
  \label{eq:QM-C05-S03-total-decay}
\end{equation}
这类非 Hermitian 势是只保留某些通道的有效描述：概率离开所保留的子空间进入吸收装置、反应产物或其他未显式建模的自由度。把环境和所有通道纳入更大的封闭系统后，总演化仍应是酉的。完整推导见
\relatedexercise{QM-C05-S03-E01}{教材练习 5.3.1：复吸收势与概率衰减}。

\textbf{一分钟回顾。}

概率连续性方程 $\partial_t\rho+\nabla\cdot\mathbf j=0$ 把整体归一化提升为局域守恒；对任意固定区域，概率变化率等于负的边界净向外通量。标准概率流由 Schr\"odinger 方程推出，既依赖振幅又依赖空间相位。定态只要求 $\nabla\cdot\mathbf j=0$，束缚边界才通常进一步迫使 $\mathbf j=0$。实波函数的零净流不等于零动量确定态；左右行波可在密度中干涉而在净流中相减。负虚势则作为有效概率汇，使保留通道中的范数衰减。

\clearpage
\section{第 5.4 节：单阶跃势与一维散射}
\label{section:QM-C05-S04}

\studymeta
  {QM-C05-S04}
  {2026-09-15}
  {Shankar, \textit{Principles of Quantum Mechanics}, Second Edition}
  {第 5 章 \S 5.4；印刷页 pp.~167--175；PDF pp.~179--187}
  {讲解、教材练习 5.4.1--5.4.3 与订正已完成}

\textbf{本节主线。}

一维散射把真实的局域波包过程与便于计算的定态平面波方法连接起来。对单阶跃势，先由分段 Schr\"odinger 方程和边界匹配求入射、反射与透射振幅，再用概率流而非单纯振幅模平方定义反射率和透射率。$E>V_0$ 时量子粒子仍可部分反射；$E<V_0$ 时半无限台阶中存在无净流的倏逝波，而有限宽势垒则允许非零隧穿。

\textbf{来源边界。}阶跃势、Gaussian 入射波包、定态散射算法、反射与透射系数及 $E<V_0$ 的讨论来自教材第 5.4 节。有限势垒的完整匹配公式、虚波数的严格解释以及传播子中经典作用量的说明，是结合教材习题与对话推导补全的内容。教材脚注把本节标为较难，并允许读者学完第 7 章后再读；本笔记仍按当前所需前置知识展开全部关键步骤。

\notetopic{QM-C05-S04-T01}{真实波包、散射通道与定态算法}

\keyidea{局域波包负责描述“何时到达、如何分裂”的动态过程；动量足够尖锐时，各波数分量近似具有相同散射系数，因而可用单能量定态解计算最终反射和透射概率。}

以下取质量 $m>0$ 恒定、实势且无吸收，并考虑高度 $V_0>0$ 的单阶跃势
\begin{equation}
  V(x)=
  \begin{cases}
    0,&x<0\quad(\text{区域 I}),\\
    V_0,&x>0\quad(\text{区域 II}).
  \end{cases}
  \label{eq:QM-C05-S04-step-potential}
\end{equation}
经典粒子从左侧入射时，$E>V_0$ 将越过台阶并减速，$E<V_0$ 则完全反射。量子计算将分别显示：前者仍有非零反射，后者的波函数仍可渗入经典禁区。

严格平面波遍布全空间，不能描述一个最初位于左侧、随后碰撞台阶的粒子。教材因此从 Gaussian 波包出发：
\begin{equation}
  \psi_I(x,0)
  =(\pi\Delta^2)^{-1/4}
  \exp[\ii k_0(x+a)]
  \exp\left[-\frac{(x+a)^2}{2\Delta^2}\right].
  \label{eq:QM-C05-S04-incident-gaussian}
\end{equation}
它满足
\begin{equation}
  \expect X=-a,
  \qquad
  \expect P=\hbar k_0,
  \qquad
  \Delta X=\frac{\Delta}{\sqrt2},
  \qquad
  \Delta P=\frac{\hbar}{\sqrt2\Delta}.
  \label{eq:QM-C05-S04-gaussian-data}
\end{equation}
取 $k_0>0$、$a\gg\Delta$ 使初态远离台阶；取 $k_0\Delta\gg1$ 则使动量分布窄峰化于 $k_0$，从而近似具有确定入射能量
\begin{equation}
  E_0=\frac{\hbar^2k_0^2}{2m}.
  \label{eq:QM-C05-S04-incident-energy}
\end{equation}
这里必须接受不确定性带来的权衡：能量越尖锐，入射波包在位置空间越宽。

\begin{figure}[htbp]
  \centering
  \begin{tikzpicture}[>=Latex,scale=0.9]
    \draw[->] (-5,0)--(5,0) node[right] {$x$};
    \draw[->] (-5,0)--(-5,1.9) node[above] {$V(x)$};
    \draw[thick] (-5,0)--(0,0)--(0,1.1)--(4.5,1.1);
    \node[above] at (3.7,1.1) {$V_0$};
    \node[below] at (-2.5,0) {区域 I};
    \node[below] at (2.5,0) {区域 II};
    \draw[->] (-5,-2.5)--(5,-2.5) node[right] {$x$};
    \draw[->] (-5,-2.5)--(-5,-0.8) node[above] {$|\psi|^2$};
    \draw[dashed] (0,-2.7)--(0,1.5);
    \draw[thick,blue!70!black,domain=-4.7:-0.7,samples=80]
      plot (\x,{-2.5+1.15*exp(-2*(\x+2.7)^2)});
    \draw[->,blue!70!black] (-2.2,-1.45)--(-3.3,-1.45);
    \node[above,blue!70!black] at (-2.7,-1.45) {反射包};
    \draw[thick,red!70!black,domain=0.6:3.4,samples=80]
      plot (\x,{-2.5+0.95*exp(-4*(\x-1.8)^2)});
    \draw[->,red!70!black] (1.5,-1.5)--(2.5,-1.5);
    \node[above,red!70!black] at (2.0,-1.5) {透射包};
  \end{tikzpicture}
  \caption{上图为势能台阶，下图示意 $E_0>V_0$ 时碰撞后分离的两个波包及其传播方向。密度曲线仅作示意，面积和高度不代表特定参数下的数值结果。}
  \label{fig:QM-C05-S04-packet-scattering}
\end{figure}

在足够晚的时刻，反射包与透射包彼此分离，定义
\begin{equation}
  R=\int\abs{\psi_R(x,t)}^2\dd x,
  \qquad
  T=\int\abs{\psi_T(x,t)}^2\dd x,
  \qquad
  R+T=1.
  \label{eq:QM-C05-S04-packet-RT}
\end{equation}
$R,T$ 是重复制备同一初态后得到两类结果的概率，并非单个粒子被分割。完整动态算法是：求定态散射本征函数，投影初始包，给每个能量分量附加时间相位，再在晚时识别两个波包。窄动量极限使散射振幅可在 $k_0$ 处提出积分，最终退化为下面的定态概率流计算。

\notetopic{QM-C05-S04-T02}{\texorpdfstring{$E>V_0$}{E > V0}：分段解、边界匹配与振幅}

\keyidea{有限势阶不含 delta 奇点，因此波函数及其一阶导数在边界连续；“只从左侧入射”则排除区域 II 中从 $+\infty$ 向左传播的分量。}

对定态能量 $E>V_0$，定义
\begin{equation}
  k_1=\frac{\sqrt{2mE}}{\hbar},
  \qquad
  k_2=\frac{\sqrt{2m(E-V_0)}}{\hbar}.
  \label{eq:QM-C05-S04-wave-numbers}
\end{equation}
两区的一般振荡解为
\begin{align}
  \psi_I(x)&=A\ee^{\ii k_1x}+B\ee^{-\ii k_1x},\\
  \psi_{II}(x)&=C\ee^{\ii k_2x}+D\ee^{-\ii k_2x}.
  \label{eq:QM-C05-S04-general-propagating-solutions}
\end{align}
$A,B,C,D$ 依次对应右行入射、左行反射、右行透射和从右侧入射的左行波。给定只有左侧入射源，故 $D=0$，但 $B$ 可由台阶反射产生。

若 $\psi$ 在边界跳跃，则 $\psi''$ 含有 delta 的导数，而有限势对应的 $(V-E)\psi$ 无法与之平衡，故 $\psi$ 连续。再把定态 Schr\"odinger 方程在 $[-\epsilon,\epsilon]$ 上积分：
\begin{equation}
  \psi'(\epsilon)-\psi'(-\epsilon)
  =\frac{2m}{\hbar^2}
    \int_{-\epsilon}^{\epsilon}(V-E)\psi\,\dd x
  \longrightarrow0.
\end{equation}
有限 $V_0$ 不产生有限面积的奇异项，因此
\begin{equation}
  \psi(0^-)=\psi(0^+),
  \qquad
  \psi'(0^-)=\psi'(0^+).
  \label{eq:QM-C05-S04-finite-step-boundary}
\end{equation}
这与 delta 势“$\psi$ 连续、$\psi'$ 跳跃”形成直接对比。代入分段解得到
\begin{equation}
  A+B=C,
  \qquad
  k_1(A-B)=k_2C.
  \label{eq:QM-C05-S04-matching-equations}
\end{equation}
由 $k_1(A-B)=k_2(A+B)$ 得 $(k_1-k_2)A=(k_1+k_2)B$，于是
\begin{equation}
  \boxed{
  r\equiv\frac BA=\frac{k_1-k_2}{k_1+k_2}},
  \qquad
  \boxed{
  t\equiv\frac CA=\frac{2k_1}{k_1+k_2}}.
  \label{eq:QM-C05-S04-amplitudes-above-step}
\end{equation}
$r,t$ 是振幅比而非概率。$V_0=0$ 时 $k_1=k_2$，于是 $r=0,t=1$，这是最直接的边界自检。

\notetopic{QM-C05-S04-T03}{概率流、速度因子与波包 Jacobian}

\keyidea{散射概率比较的是单位时间通过截面的概率，而不是局部密度。平面波流满足 $j=\rho v$；两侧速度不同时，透射率必须包含速度比。}

右行和左行平面波的概率流分别为
\begin{equation}
  \psi=F\ee^{\pm\ii kx}
  \quad\Longrightarrow\quad
  j=\pm\frac{\hbar k}{m}\abs F^2.
  \label{eq:QM-C05-S04-plane-wave-current}
\end{equation}
因此
\begin{equation}
  j_{\mathrm{in}}=\frac{\hbar k_1}{m}\abs A^2,
  \qquad
  j_{\mathrm{ref}}=-\frac{\hbar k_1}{m}\abs B^2,
  \qquad
  j_{\mathrm{tr}}=\frac{\hbar k_2}{m}\abs C^2.
  \label{eq:QM-C05-S04-three-currents}
\end{equation}
区域 I 中左右行波的交叉项在概率流的反对称组合中抵消，详见式~\eqref{eq:QM-C05-S03-counterpropagating-current}。按流的正大小定义
\begin{align}
  R&=\frac{\abs{j_{\mathrm{ref}}}}{j_{\mathrm{in}}}
   =\abs{\frac BA}^2
   =\boxed{\left(\frac{k_1-k_2}{k_1+k_2}\right)^2},
  \label{eq:QM-C05-S04-reflection-coefficient}\\
  T&=\frac{j_{\mathrm{tr}}}{j_{\mathrm{in}}}
   =\frac{k_2}{k_1}\abs{\frac CA}^2
   =\boxed{\frac{4k_1k_2}{(k_1+k_2)^2}}.
  \label{eq:QM-C05-S04-transmission-coefficient}
\end{align}
所以
\begin{equation}
  R+T=1.
  \label{eq:QM-C05-S04-current-conservation}
\end{equation}
只有入射区与透射区波数相同，才可把 $T$ 简写为 $\abs{C/A}^2$。例如 $E\to V_0^+$ 时 $k_2\to0$，虽有 $C/A\to2$，但透射速度趋零，故 $T\to0$。

教材练习 5.4.1 从完整波包演化验证同一速度因子。令
\begin{equation}
  k_2(k_1)=\sqrt{k_1^2-\frac{2mV_0}{\hbar^2}},
  \qquad
  k_{20}=k_2(k_0),
  \qquad
  q=k_1-k_0.
\end{equation}
要求入射谱远离传播阈值，即 $k_0-\sqrt{2mV_0}/\hbar\gg1/\Delta$，才能在整个有效窄峰上作一阶展开：
\begin{equation}
  \frac{\dd k_2}{\dd k_1}=\frac{k_1}{k_2},
  \qquad
  k_2(k_1)\simeq k_{20}+\frac{k_0}{k_{20}}q.
  \label{eq:QM-C05-S04-k2-linearization}
\end{equation}
透射包络因而依赖
\begin{equation}
  X=\frac{k_0}{k_{20}}x+a-\frac{\hbar k_0}{m}t.
  \label{eq:QM-C05-S04-transmitted-envelope-coordinate}
\end{equation}
其中心和群速度为
\begin{equation}
  x_c(t)=\frac{\hbar k_{20}}{m}t-\frac{k_{20}}{k_0}a,
  \qquad
  \dot x_c=\frac{\hbar k_{20}}m.
  \label{eq:QM-C05-S04-transmitted-center}
\end{equation}
变量替换给出
\begin{equation}
  \dd x=\frac{k_{20}}{k_0}\dd X,
  \label{eq:QM-C05-S04-packet-jacobian}
\end{equation}
于是归一化包络的总透射概率为
\begin{equation}
  \boxed{
  P_T\simeq\frac{k_{20}}{k_0}\abs{\frac CA}^2}.
  \label{eq:QM-C05-S04-packet-transmission}
\end{equation}
这正是 $v_2/v_1$，振幅比取峰值 $k_0$ 处。上述形状与中心采用窄谱、线性相位近似；若计算长时间扩散，须一致保留空间相位 $k_2(k_1)x$ 和时间相位中的二阶项，不能把线性近似无限外推。总透射概率还可直接由动量积分的变量替换得到，无需近似波包形状。完整推导见
\relatedexercise{QM-C05-S04-E01}{教材练习 5.4.1：波包与定态透射率的一致性}。

\notetopic{QM-C05-S04-T04}{\texorpdfstring{$E<V_0$}{E < V0}：倏逝波、全反射与有限势垒隧穿}

\keyidea{半无限台阶中的倏逝波具有非零概率密度但零净流；有限宽势垒提供第二个边界，使衰减后的振幅重新接入传播通道，从而产生指数小但非零的透射。}

当 $0<E<V_0$ 时，区域 I 仍取
\begin{equation}
  \psi_I=A\ee^{\ii kx}+B\ee^{-\ii kx},
  \qquad
  k=\frac{\sqrt{2mE}}{\hbar},
\end{equation}
而区域 II 满足 $\psi''=\kappa^2\psi$，其中
\begin{equation}
  \kappa=\frac{\sqrt{2m(V_0-E)}}{\hbar}.
  \label{eq:QM-C05-S04-kappa}
\end{equation}
排除 $x\to+\infty$ 发散项后，
\begin{equation}
  \psi_{II}=C\ee^{-\kappa x}.
  \label{eq:QM-C05-S04-evanescent-wave}
\end{equation}
边界匹配为 $A+B=C$ 和 $\ii k(A-B)=-\kappa C$，解得
\begin{equation}
  r=\frac BA=\frac{k-\ii\kappa}{k+\ii\kappa},
  \qquad
  t=\frac CA=\frac{2k}{k+\ii\kappa}.
  \label{eq:QM-C05-S04-below-step-amplitudes}
\end{equation}
分子和分母模长相同，因此
\begin{equation}
  \boxed{R=\abs r^2=1}.
  \label{eq:QM-C05-S04-total-reflection}
\end{equation}
$r$ 一般仍有非平凡相位，故全反射不等于波函数原样返回。区域 II 的密度与概率流分别为
\begin{equation}
  \rho_{II}=\abs C^2\ee^{-2\kappa x}\neq0,
  \qquad
  \boxed{j_{II}=0}.
  \label{eq:QM-C05-S04-evanescent-density-current}
\end{equation}
由于 $\psi'_{II}=-\kappa\psi_{II}$，流公式中两项均为 $-\kappa|\psi_{II}|^2$，相减为零。所以禁区附近仍可探测到粒子，但没有向 $+\infty$ 输运的概率。这里的 $j_{II}=0$ 是单能量定态结论；真实波包碰撞期间可以暂时流入禁区，再流回左侧，不代表禁区概率从不变化。

倏逝波不是“虚概率波”，其密度仍是实数且非负。形式记号 $k_2=\ii\kappa$ 只表示振荡解变为指数解；$\ii\hbar\kappa$ 不是自伴动量算符的可测量本征值。

若高势只存在于 $\abs x<a$，禁区宽度 $L=2a$，则右侧重新具有传播通道。势垒内部必须保留 $\psi=F\ee^{\kappa x}+G\ee^{-\kappa x}$；两个指数项都在有限区间内有界，其交叉项可以携带非零流：
\begin{equation}
  j=\frac{2\hbar\kappa}{m}\operatorname{Im}(FG^*).
\end{equation}
因此隧穿过程的流穿过整个势垒，并非只在第二边界才重新出现。对 $0<E<V_0$，完整匹配结果为
\begin{equation}
  \boxed{
  T=
  \left[
    1+\frac{V_0^2}{4E(V_0-E)}
    \sinh^2(2\kappa a)
  \right]^{-1}},
  \qquad
  R=1-T.
  \label{eq:QM-C05-S04-finite-barrier-transmission}
\end{equation}
厚势垒 $\kappa a\gg1$ 时
\begin{equation}
  T\simeq
  \frac{16E(V_0-E)}{V_0^2}
  \ee^{-4\kappa a}
  =\frac{16E(V_0-E)}{V_0^2}
  \ee^{-2\kappa L}.
  \label{eq:QM-C05-S04-thick-barrier}
\end{equation}
因此增大势垒宽度、高度差 $V_0-E$ 或质量都会指数压低隧穿率，但有限参数下不把它严格变为零。完整匹配见
\relatedexercise{QM-C05-S04-E02B}{教材练习 5.4.2(b)：有限方势垒}；排斥 delta 势的散射见
\relatedexercise{QM-C05-S04-E02A}{教材练习 5.4.2(a)：delta 势散射}。

对更一般、在 $x\to\pm\infty$ 足够快趋于常数的实势，当能量高于两侧渐近势能时，仍可指定左侧入射通道，求相应的渐近平面波，再以概率流之比定义 $R,T$。理想平面波不可普通归一化，但其流之比有限；它是窄动量波包极限下的稳态计算工具。

\textbf{一分钟回顾。}

散射的真实对象是局域波包，定态平面波则把窄能量分布的动态问题化为边界匹配。有限势阶要求 $\psi$ 与 $\psi'$ 连续；只从左侧入射则排除右侧来波。反射率和透射率必须比较概率流，故透射振幅模平方还要乘速度比。$E<V_0$ 的半无限台阶允许倏逝密度却没有透射流；有限势垒的第二边界把倏逝振幅重新接入传播通道，产生 $T\propto\ee^{-2\kappa L}$ 的隧穿。波函数非零、概率流非零和穿过整个势垒是三个不同判断。

\clearpage
\section{第 5.5 节：双缝实验}
\label{section:QM-C05-S05}

\studymeta
  {QM-C05-S05}
  {2026-09-16}
  {Shankar, \textit{Principles of Quantum Mechanics}, Second Edition}
  {第 5 章 \S 5.5；印刷页 p.~175 末至 p.~176 开头；PDF pp.~187--188}
  {讲解与一道自编短应用题已完成并核对}

\textbf{来源边界。}教材用一小段文字，将第三章双缝实验重新解释为波函数的传播、干涉及 Born 概率规则的应用，没有编号习题。下文的 Helmholtz 方程、几何展开和通量说明是对话中的推导补充；条纹间距联系教材第 3.2 节及图 3.1（印刷页 pp.~108--109，PDF pp.~121--122）。

\notetopic{QM-C05-S05-T01}{从边界条件到干涉图样与概率解释}

\keyidea{经典波动理论提供传播和干涉的数学方法；量子力学再用 Born 规则解释波函数模平方。相同的波动方程不意味着波函数与经典物理场具有相同含义。}

\textbf{装置与波动方程。}
考虑非相对论、有质量的粒子，静止挡板的实体部分以无限高势垒理想化，波函数在硬壁边界为零；狭缝是开放通道，不能也设为零。实际入射态应是波包，动量分布足够窄时可用平面波近似计算图样。对自由区域中的定态，
\begin{equation}
  \psi(\mathbf r,t)=\phi(\mathbf r)\ee^{-\ii Et/\hbar},
  \qquad
  -\frac{\hbar^2}{2m}\nabla^2\phi=E\phi
  \quad\Longleftrightarrow\quad
  \nabla^2\phi+k^2\phi=0,
  \label{eq:QM-C05-S05-helmholtz}
\end{equation}
其中 $k=\sqrt{2mE}/\hbar=p/\hbar$、$\lambda=2\pi/k=h/p$。这与经典单色波的 Helmholtz 方程同形，故可借用波动理论处理开口后的衍射和传播。

\begin{figure}[htbp]
  \centering
  \begin{tikzpicture}[>=Latex,scale=0.85]
    \draw[very thick] (0,-1.55)--(0,-0.87);
    \draw[very thick] (0,-0.63)--(0,0.63);
    \draw[very thick] (0,0.87)--(0,1.55);
    \draw[very thick] (5,-1.55)--(5,1.8);
    \draw[dashed] (-3.1,0)--(5,0);
    \draw[->] (-3.1,0.4)--(-1.9,0.4);
    \node[above] at (-2.5,0.4) {入射};
    \node[above] at (0,1.55) {双缝挡板};
    \node[above] at (5,1.8) {探测屏};
    \node[left] at (0,0.75) {$S_1$};
    \node[left] at (0,-0.75) {$S_2$};
    \draw[<->] (-0.8,-0.75)--node[left] {$a$} (-0.8,0.75);
    \draw[blue!70!black] (0,0.75)--node[above] {$r_1$} (5,1.15);
    \draw[red!70!black] (0,-0.75)--node[below] {$r_2$} (5,1.15);
    \fill (5,1.15) circle (2pt);
    \draw[<->] (5.4,0)--node[right] {$x$} (5.4,1.15);
    \draw[<->] (0,-1.9)--node[below] {$d$} (5,-1.9);
  \end{tikzpicture}
  \caption{双缝几何示意（按教材图 3.1 自绘，不按比例）。$a$ 为缝距，$d$ 为挡板到探测屏的距离，$x$ 为屏上横向坐标。}
  \label{fig:QM-C05-S05-geometry}
\end{figure}

\textbf{相干叠加与条纹间距。}
在教材的开口近似中，两缝是同一入射波的等效源，并非创造两个新粒子。写总场为 $\phi=\phi_1+\phi_2$，两项指同一双缝装置下的相干贡献；严格计算须解整个边界问题，不能一般地把改变边界后两个单缝实验的解直接相加。忽略两路振幅大小的差异、设两缝初相位相同，在观察位置处取
\begin{equation}
  \phi_1=A\ee^{\ii kr_1},\qquad
  \phi_2=A\ee^{\ii kr_2},\qquad
  \delta=k(r_2-r_1).
\end{equation}
则
\begin{align}
  \rho=\abs{\phi_1+\phi_2}^2
  &=\abs{\phi_1}^2+\abs{\phi_2}^2
    +\phi_1^*\phi_2+\phi_2^*\phi_1\\
  &=2\abs A^2(1+\cos\delta)
   =\boxed{4\abs A^2\cos^2(\delta/2)}.
  \label{eq:QM-C05-S05-interference}
\end{align}
路程差为 $n\lambda$ 时相长，为 $(n+1/2)\lambda$ 时相消；完全相消还要求两路振幅大小相同。几何上
\begin{equation}
  r_1=\sqrt{d^2+(x-a/2)^2},\qquad
  r_2=\sqrt{d^2+(x+a/2)^2}.
\end{equation}
在远场近轴条件 $a,\abs x\ll d$ 下，使用 $\sqrt{d^2+u^2}\simeq d+u^2/(2d)$ 得
\begin{equation}
  r_2-r_1\simeq\frac{(x+a/2)^2-(x-a/2)^2}{2d}
  =\frac{ax}{d},\qquad
  \boxed{\Delta x\simeq\frac{\lambda d}{a}=\frac{hd}{pa}}.
  \label{eq:QM-C05-S05-fringe-spacing}
\end{equation}
相邻亮纹的相位差增加 $2\pi$，所以间距随动量增大而减小。有限缝宽会带来衍射包络；$A$ 是局部近似振幅，不能视作整张屏上的常数。相位与动量的短应用见\relatedexercise{QM-C05-S05-E01}{自编应用题}。

\textbf{概率密度、计数与单次事件。}
Born 规则给出可归一化波包在有限区域 $\Omega$ 内的概率
\begin{equation}
  P_\Omega(t)=\int_\Omega\abs{\psi(\mathbf r,t)}^2\dd^3r.
\end{equation}
不能把某点的密度值直接当作单点概率。单次探测产生局域事件，重复实验的统计分布显示干涉，不需要不同粒子相互碰撞，也不意味着一个粒子被分成两份。

探测屏的到达计数更精确地涉及概率通量；在理想单向探测、近轴且纵向速度近似相同的条件下，$j_n\simeq v_n\abs\psi^2$，计数图样才与相应的密度图样成正比。教材在这里采用这一理想化解释。边界条件与 Schr\"odinger 方程决定振幅，而概率含义来自 Born 公设，并不是由经典波动理论推导出的新结论。
